%%%%%%%%%%  scan idx 290  |  folio 277  %%%%%%%%%%
Il
résulte aussitôt de (1.4) que \(\varphi\) est inverse de l'homomorphisme canonique \(k(A)\longrightarrow k(D_A^b(C))\).

\textbf{Remarque 1.6.1.} Dans la situation de (1.6), le foncteur canonique \(D^b(A)\longrightarrow D_A^b(C)\) n'est pas en général une équivalence de catégories, comme le montre l'exemple où \(S\) est un schéma, \(C\) la catégorie des \(\mathcal O_S\)-Modules, \(A\) la catégorie des \(\mathcal O_S\)-Modules quasi-cohérents (II Ap. I 0.3). Voir cependant le critère (II 2.1.3).

\textbf{Définition 1.7.} Soient \(C\) une catégorie triangulée (resp. abélienne), et \(C'\) une partie de \(\Ob C\). Nous dirons que \(C'\) est un \emph{sous-ensemble exact} (cf. EGA III 3.1.1) si, pour tout triangle distingué \(E'\to E\to E''\to E'[1]\) (resp. toute suite exacte \(0\to E'\to E\to E''\to0\)) tel (resp. telle) que deux des termes \(E',E,E''\) appartiennent à \(C'\), le troisième y appartient également. Nous dirons que \(C'\) est \emph{dense} dans \(\Ob C\) si \(\Ob C\) est le plus petit sous-ensemble exact de \(\Ob C\) contenant \(C'\).

\textbf{Lemme 1.8.} Soient \(C\) une catégorie triangulée (resp. abélienne) et \(C'\) une partie dense de \(\Ob C\). Alors les éléments \(\cl_C(E)\), pour \(E\in C'\), engendrent \(k(C)\).

\emph{Preuve.} Posons \(C'=C_0\), et, pour \(i\geq1\), définissons \(C_i\) par récurrence comme l'ensemble des objets de \(C\) figurant dans un triangle distingué (resp. une suite exacte courte) dont deux des termes appartiennent à \(C_{i-1}\). Les \(C_i\) forment une filtration croissante de \(\Ob C\), qui, par hypothèse, est exhaustive. Pour \(i\geq0\), notons \(k(C)_i\) le sous-groupe de \(k(C)\) engendré par les éléments \(\cl_C(E)\) pour \(E\in C_i\). Il est immédiat sur les définitions que l'on a, pour tout \(i\geq0\), \(k(C)_i=k(C)_{i+1}\), donc \(k(C)_i=k(C)_0\). Soit alors \(x\in k(C)\); 
